\documentclass[11pt]{amsart}
\usepackage[T1]{fontenc}
\usepackage{lmodern}
\usepackage{amsmath,amssymb,amsthm}
\usepackage[margin=1.1in]{geometry}
\usepackage{booktabs}
\usepackage{hyperref}

\newtheorem{theorem}{Theorem}[section]
\newtheorem{lemma}[theorem]{Lemma}
\newtheorem{proposition}[theorem]{Proposition}
\newtheorem{corollary}[theorem]{Corollary}
\newtheorem{conjecture}[theorem]{Conjecture}
\newtheorem{question}[theorem]{Question}
\theoremstyle{definition}
\newtheorem{definition}[theorem]{Definition}
\newtheorem{remark}[theorem]{Remark}
\newtheorem{example}[theorem]{Example}

\newcommand{\cD}{\mathcal D}
\newcommand{\lean}[1]{\quad\emph{(Lean: \texttt{#1}.)}}
\emergencystretch=3em

\title[Fan's lemma at equality]{Fan's lemma at equality:\\ one residual configuration behind four
Lov\'asz-type problems, a protected-blank form of Fan's lemma, and why flowers do not apply}
\author{Justin Williams}
\date{November 2026 (working draft)}

\begin{document}
\begin{abstract}
Four lines of work on Gallai's path-decomposition conjecture stop at what looks like the same
place: the one-sided avoidance conjecture $L_1^*$ in its blocked case, the locality statement
$L_\infty$ ($Q_c\subseteq N[c]$ in all-odd graphs), the end-edge conjecture EE when $u$ and the even
vertex $z$ have two common neighbours, and the base step of the walk theorem for the
terminal-matching property PTM$_2$. All four ask for a Lov\'asz decomposition with a one-sided
property of one protected path, and all four use Lov\'asz's induction with Fan's lemma. We make the
common obstruction precise.

\begin{itemize}
\item \emph{The configuration.} On a Lov\'asz decomposition, Fan's chain is forced
(Proposition~\ref{prop:forced}). The obstruction is a Lov\'asz Case~I step $t\to w$ whose target $w$
is adjacent to an end $p$ of the protected path. In the leaf graph, $p$ is then a blank of $w$, and
Fan's inequality is the equality $b(t)=1=|Z(w)|$. The invariant survives exactly when the forced
chain avoids $p$. Exhaustive data confirm this dichotomy without exception in the $L_\infty$
residual and in the blocked residual of $L_1^*$, and they show it is \emph{literally} the obstruction
in three of the four settings: $L_\infty$ (with $p$ the far end of the leaf path), the minimal
blocked shape of $L_1^*$ (with $p=c$), and PTM instances without an admissible walk (with $p$ the
short end). The blocked case of $L_1^*$ has in addition a hypothesis-side failure: the instance on
$G-gu$ is good but not admissible. The $K_4$/$W_4$ base failures of PTM are of a third kind.
\item \emph{A protected-blank form of Fan's lemma} (Theorem~\ref{thm:protected}): a blank that no
chain vertex reaches as a predecessor does not count against the supply. As a consequence, the
tight transfer of $L_1^*$ goes through under an explicit condition on the decomposition
(Theorem~\ref{thm:l1tight}); at order $7$ the condition can be met in every tight instance.
\item \emph{Unification.} $L_1^*$ on blocked instances implies both EE (known) and $L_\infty$
(Theorem~\ref{thm:linfty}). So three of the four residuals are one statement.
\item \emph{Flowers do not apply} (Section~\ref{sec:flower}). The flower move and receipt-or-flower
always produce a vertex with two ends, and the equality that forces a flower cannot occur with
$|T|\ge2$ when every vertex has at most one end. No flower move ends in a Lov\'asz decomposition,
so none of the four residuals can be closed by it.
\item \emph{Test~1 at degree four} (Theorem~\ref{thm:test1}): the first deletion test of the
spine holds at every vertex $z$ of degree four with $J-z$ connected, by two-step descent, without
PTM.
\end{itemize}
Everything stated as a theorem is formalised in Lean~4 with Mathlib
(\texttt{lean/GallaiLean/TightFan.lean}, no \texttt{sorry}, standard axioms only), except
Theorem~\ref{thm:tightwalk}, which is proved here on paper. Computations are in a separate
section; no proof depends on them.
\end{abstract}
\maketitle

\section{Introduction}\label{sec:intro}

All graphs are finite and simple. A \emph{path decomposition} $\cD$ of $G$ partitions $E(G)$ into
paths with at least one edge; $b_\cD(v)$ is the number of paths of $\cD$ with an end at $v$, and
$b_\cD(v)\equiv d(v)\pmod2$. A vertex is \emph{even} if its degree is positive and even. A
\emph{Lov\'asz decomposition} has $b(v)=d(v)\bmod2$ for every $v$; Lov\'asz~\cite{L68} proved that
a graph with at most one even vertex has one. For an odd vertex $c$, $Q_c$ is the path ending at
$c$. Write $Z_\cD(w)=\{v\in N(w):b_\cD(v)=0\}$ for the \emph{blanks} of $w$.

Lov\'asz's induction deletes an edge $gu$ at the even vertex $g$, decomposes $G-gu$ (where $u$ is
the even vertex), and restores $gu$ by Fan's lemma~\cite{Fan05} with source $g$ and target $u$.
Fan's lemma keeps every vertex set except that one path gains $u$~\cite[Lemma~2.2]{RD}. So a
property of the vertex set of one protected path survives if the induction supplies it on $G-gu$ and
the transfer does not spoil it. Four statements of the project are proved this way except in one
residual configuration each:
\begin{enumerate}
\item[(R1)] $L_\infty$ \cite[\S8]{RD}: in a graph with no even vertex, some Lov\'asz decomposition
has $V(Q_c)\subseteq N[c]$. Through a leaf $\ell$ at $c$, Lov\'asz's induction proves it except
when the even vertex $z$ is at distance $3$ from $c$, every neighbour of $z$ is at distance $2$,
and $d(z)\ge4$ or $z$ lies on a triangle.
\item[(R2)] $L_1^*$ \cite{L1S}: with at most one even vertex $g$, some Lov\'asz decomposition has
$Q_c\cap A=\emptyset$ whenever $A\cap(N[c]\cup N(g))=\emptyset$. It is reduced to \emph{blocked}
instances: $c\notin N(g)$ and every neighbour of $g$ is adjacent to $A\setminus\{g\}$.
\item[(R3)] EE \cite{PE}: if $z$ is the only even vertex and $u\in N(z)$, then $uz$ is terminal
at $u$ in some Lov\'asz decomposition. It is open when $|N(u)\cap N(z)|\ge2$, and it follows from
the blocked case of $L_1^*$ \cite[Corollary~4.4]{L1S}.
\item[(R4)] PTM$_2$ \cite{TM}: two prescribed terminal edges $ab$, $xy$ on different paths. The
walk theorem reduces it to a separation problem in a base graph; that base problem fails at local
copies of the $K_4$/$W_4$ obstruction, and in instances with no admissible walk.
\end{enumerate}
It had been suggested that all four are ``the same tight case of Fan's inequality'' and that the
flower move of~\cite[\S7]{ES}, which is designed for Fan's equality case, might close them. This
paper tests both suggestions.

\subsection*{Results}
\begin{enumerate}
\item Section~\ref{sec:fan}: Fan's chain is forced on Lov\'asz decompositions
(Proposition~\ref{prop:forced}), and Fan's lemma holds with \emph{protected blanks}
(Theorem~\ref{thm:protected}).
\item Section~\ref{sec:four}: the precise configuration in each of (R1)--(R4), the dichotomy
``the invariant survives iff the forced chain avoids $p$'', the tight transfer for $L_1^*$
(Theorem~\ref{thm:l1tight}), $L_\infty$ from blocked $L_1^*$ (Theorem~\ref{thm:linfty}), and a
tight walk theorem for PTM (Theorem~\ref{thm:tightwalk}). The answer to ``literally one
configuration?'' is: (R1), the minimal shape of (R2), and the no-walk part of (R4) are literally one
configuration; (R2) also has a hypothesis-side part; (R3) is (R2); the $K_4$/$W_4$ part of (R4) is
different.
\item Section~\ref{sec:flower}: the flower move never ends in a Lov\'asz decomposition.
\item Section~\ref{sec:test1}: Test~1 of the spine at degree four, without PTM.
\item Section~\ref{sec:comp}: the computations (all exhaustive unless stated).
\item Section~\ref{sec:open}: what is left, and the precise new obstruction.
\end{enumerate}

\subsection*{Honest status}
None of (R1)--(R4) is closed here. What is proved is the exact form of the obstruction, a
conditional closure of the tight case, the reduction of (R1) and (R3) to one statement, and a
complete negative answer for flower moves. The computations show that the remaining difficulty is
in the \emph{induction hypothesis}, not in a missing move: in the $L_\infty$ residual there are
instances in which every target admits an induction decomposition that fails
(Section~\ref{sec:comp}, E26(B)).

\section{Fan's chain, and Fan's lemma with protected blanks}\label{sec:fan}

Let $\cD$ decompose $K$, let $w$ be a vertex and $t\notin N[w]$. For a path $P\in\cD$ with an end
$x$ that contains $w$, the \emph{predecessor} of $w$ read from $x$ is the vertex of $P$ just before
$w$ when $P$ is read from $x$. A \emph{chain} from $t$ is a sequence $t=x_0,x_1,\dots,x_k$ of
distinct vertices in which $x_{i+1}$ is a predecessor of $w$ read from $x_i$, and some path with an
end at $x_k$ avoids $w$. Executing a chain (re-route at $x_0,\dots,x_{k-1}$, then append $x_kw$)
gives a decomposition of $K+tw$ of the same size with $b-\mathbf 1_t+\mathbf 1_w$, every path of
which lies in a path of $\cD$ together with $w$ \cite[Lemma~2.2]{RD}\lean{fan\_exec\_supp}. Fan's
lemma says a chain exists when $\sum_{t\in T}b(t)>|Z(w)|$ for a set $T$ of non-neighbours of
$w$\lean{fan\_exists\_chain}.

\begin{proposition}[the chain is forced]\label{prop:forced}
If $b_\cD(x)\le1$, the predecessor of $w$ read from $x$ is unique when it exists. Hence, if
$b_\cD\le1$ everywhere, Fan's chain from $t$ is determined by $\cD$, $t$ and $w$: $x_{i+1}$ is the
predecessor of $w$ on the unique path ending at $x_i$.\lean{fanPred\_unique}
\end{proposition}

\begin{proof}
At most one path has an end at $x$, and it has only one end there. In a path, the vertex before
$w$ (in a fixed reading direction) is unique, because the support has no repeated vertex.
\end{proof}

So in the Case~I steps of all four problems (where $\cD$ is a Lov\'asz decomposition of $G-gu$) the
only freedom is the choice of the target $u$ and of the decomposition supplied by the induction.

\begin{theorem}[Fan's lemma with protected blanks]\label{thm:protected}
Let $\cD$ decompose $K$, let $T$ be a set of non-neighbours of $w$ with $w\notin T$, and let $B$ be
a set of \emph{protected} vertices with $B\cap T=\emptyset$, such that no vertex of
$T\cup N(w)$ has a vertex of $B$ as its predecessor of $w$. If
\[\bigl|\{v\in N(w): b_\cD(v)=0,\ v\notin B\}\bigr|<\sum_{t\in T}b_\cD(t),\]
then some $t\in T$ starts a chain, and the chain avoids $B$. Consequently $K+tw$ has a
decomposition of the same size, with end counts $b_\cD-\mathbf 1_t+\mathbf 1_w$ (in particular
unchanged on $B$), every path of which lies in a path of $\cD$ together with $w$.
\lean{fan\_exists\_chain\_protected, fan\_grouped\_protected\_supp, fan\_single\_protected\_supp}
\end{theorem}

\begin{proof}
Suppose no $t\in T$ starts a chain, and let $R$ be the set of last vertices of the sequences
$t=x_0,\dots,x_j$ ($t\in T$) of distinct vertices with each $x_{i+1}$ a predecessor read from $x_i$.
Every vertex of $R\setminus T$ is a predecessor, hence adjacent to $w$. By induction along the
sequence, every vertex of $R$ lies in $T\cup N(w)$, so its successor in the sequence is not in $B$;
thus $R\cap B=\emptyset$. Since there is no chain, every path with an end in $R$ contains $w$, and
the predecessor read from that end is in $R$. Map each end in $R$ to the edge between $w$ and that
predecessor. Two ends of one path give the two different edges of the path at $w$, and different
paths have disjoint edges. So $\sum_{x\in R}b(x)\le|N(w)\cap R|$. On the other hand $T$ and
$N(w)\cap R$ are disjoint subsets of $R$, and every vertex of $N(w)\cap R$ with no end is a blank
outside $B$. Hence
\[\sum_{t\in T}b(t)+|N(w)\cap R|-\bigl|\{v\in N(w): b(v)=0,\ v\notin B\}\bigr|\le\sum_{x\in R}b(x)
\le |N(w)\cap R|,\]
which contradicts the hypothesis. The chain found avoids $B$ by the same induction. The rest is
Lemma~2.2 of~\cite{RD}.
\end{proof}

With $B=\emptyset$ this is Fan's lemma. The case that matters is $B=\{p\}$ for one blank $p$ with
$\sum_Tb=1$: Fan's inequality $1>|Z(w)|=1$ fails, but the protected form still applies.

\begin{corollary}[one protected blank]\label{cor:oneblank}
Let $b_\cD\le1$, $b_\cD(t)=1$, and $Z(w)=\{p\}$ with $p\ne t$. Let $P_p$ be the path through the
edge $pw$ and $e$ its end on the side of $p$ (reading from $e$, the vertex before $w$ is $p$). Then:
\begin{enumerate}
\item the forced chain from $t$ either ends at a path avoiding $w$, or reaches $p$ and stops
there; it reaches $p$ exactly when it reaches $e$;
\item if $e\notin\{t\}\cup N(w)$, the conclusion of Theorem~\ref{thm:protected} holds with $B=\{p\}$.
\end{enumerate}
\end{corollary}

\begin{proof}
$p$ has no end, so it is interior to $P_p$, and $pw$ is its only edge to $w$. Reading $P_p$ from
its other end, the vertex before $w$ is $w$'s other neighbour on $P_p$, not $p$. So $e$ is the only
vertex having $p$ as predecessor. Chain vertices after $t$ are adjacent to $w$ and have one end
(Proposition~\ref{prop:forced}); the chain is injective (two vertices with the same predecessor
would be the two ends of one path, and then the predecessors differ), so it terminates, and it
cannot continue from $p$. This gives (1). For (2), the hypothesis of Theorem~\ref{thm:protected}
on $B=\{p\}$ is exactly $e\notin T\cup N(w)$.
\end{proof}

\begin{definition}[the tight configuration]\label{def:tight}
A \emph{tight transfer} is a Case~I step: a decomposition $\cD'$ of $K'=K-tw$ with $b\le1$,
$b(t)=1$, $b(w)=0$, $Z(w)=\emptyset$, and a protected path $P^*\in\cD'$ with an end $p\in N(w)$,
where the property to be kept is a property of $V(P^*)$ that the final append of $w$ to $P^*$
(or the re-routing of $P^*$'s end away from $p$) would destroy. In the leaf graph that freezes $P^*$,
$p$ is a blank of $w$, and Fan's inequality reads $b(t)=1=|Z(w)|$.
\end{definition}

\section{The four residuals, precisely}\label{sec:four}

\subsection{\texorpdfstring{$L_\infty$}{L-infinity}}
Let $K$ have a leaf $\ell$ with neighbour $c$, at most one even vertex $z$, and put
$B=\{\ell\}\cup N[c]$; the statement GL is that some Lov\'asz decomposition has its $\ell$-path inside
$B$ \cite[\S8]{RD}. In the residual configuration, delete $zu$ for some $u\in N(z)$ (so $u\notin B$),
take a decomposition $\cD'$ of $K-zu$ whose $\ell$-path $P^*$ lies in $B$, and let $y$ be the far end
of $P^*$.

\begin{proposition}\label{prop:linfty}
The Fan transfer $z\to u$ applied to $\cD'$ keeps $P^*$ unchanged unless the forced chain from $z$
reaches $y$; if it reaches $y$, then $y\in N(u)$, the chain stops at $y$, and $P^*$ gains $u\notin
B$. If $e$, the end on $y$'s side of the path through $yu$, is not in $\{z\}\cup N(u)$, the chain
does not reach $y$.
\end{proposition}

\begin{proof}
$P^*$ does not contain $u$, so it is never re-routed, and it changes only if the chain ends at one of
its ends. The leaf $\ell$ is not adjacent to $u$ and is not $z$, so it is never a chain vertex. If
the chain reaches $y$, the only path ending at $y$ is $P^*$, which avoids $u$, so $u$ is appended to
$P^*$. The last sentence is Corollary~\ref{cor:oneblank} in the graph with a second leaf hung at $y$
(which makes $y$ a blank of $u$).
\end{proof}

So the $L_\infty$ residual is the tight configuration with $p=y$. In the data (E26(B), all $133$
residual instances of order $9$, all $409{,}416$ pairs $(u,\cD')$), the transfer fails exactly when
the forced chain reaches $y$.

\subsection{\texorpdfstring{$L_1^*$}{L1*}: the blocked case}
For a blocked instance $(G,c,A)$ and $u\in N(g)$, the instance $(G-gu,c,A)$ is never admissible,
because $u$ has a neighbour in $A\setminus\{g\}$. This is a \emph{hypothesis-side} failure. The data
show that it is the only failure when a good target lies outside $N(c)$: at order $9$, every one of
the $2{,}414$ irreducible blocked instances with $|A|\le2$ has a target $u$ for which $(G-gu,c,A)$
is good although inadmissible. When such $u\notin N(c)$ exists, the protected transfer of
\cite[Lemma~2.4]{L1S} finishes the instance. In the others, including \emph{every} instance of the
minimal shape ($g\in A$, $d(g)=2$, $N(g)\subseteq N(c)$ independent), every good target lies in
$N(c)$. Then, with a leaf hung at $c$, $c$ is a blank of $u$, and the transfer is tight with $p=c$:
the plain transfer succeeds exactly when the forced chain avoids $c$ (E26(A)).

\begin{theorem}[the tight transfer for $L_1^*$]\label{thm:l1tight}
Let $g$ be the only even vertex of $G$, let $c$ be odd, and let $x\in N(g)$ with $x\ne c$ and
$x\notin A$. Let $\cD_3$ decompose the leaf graph $G^{+c}-gx$ (a leaf $\ell$ hung at $c$,
the edge $gx$ removed) with $b(c)=0$, $b(\ell)=1$ and the Lov\'asz end counts of $G-gx$ elsewhere,
and suppose that the leaf path avoids $A$ and that no vertex of $\{g\}\cup N_{G-gx}(x)$ has $c$ as
its predecessor of $x$ in $\cD_3$. Then $G$ has a Lov\'asz decomposition with
$Q_c\cap A=\emptyset$.\lean{l1\_tight\_transfer}
\end{theorem}

\begin{proof}
The blanks of $x$ in the leaf graph are $c$ (if $c\in N(x)$) and nothing else: every other neighbour
of $x$ in $G-gx$ is odd and keeps its end, and $\ell$ is not adjacent to $x$. Apply
Theorem~\ref{thm:protected} with $T=\{g\}$, $B=\{c\}$: $b(g)=1>0$. The new decomposition of $G^{+c}$
has $b(g)=0$, $b(x)=1$, $b(c)=0$, and its leaf path lies in the old leaf path plus $x\notin A$.
Delete the leaf \cite[Lemma~2.3(ii)]{L1S}\lean{avoid\_of\_apexSep}.
\end{proof}

The hypothesis on $\cD_3$ holds, for instance, when the path of $\cD'$ through $cx$ is $Q_c$ with
terminal edge $xc$ (the case already used in \cite{L1S}), or when its end on $c$'s side lies outside
$\{g\}\cup N(x)$. At order $7$ every tight blocked instance has a good decomposition of $G-gx$
satisfying it, and the condition never fails to deliver (E26(A)).

\subsection{EE, and \texorpdfstring{$L_\infty$}{L-infinity} from blocked \texorpdfstring{$L_1^*$}{L1*}}
\begin{theorem}\label{thm:linfty}
If every blocked instance of $L_1^*$ is good, then $L_\infty$ and EE hold.
\lean{lInfty\_of\_l1Blocked, ee\_of\_l1Blocked}
\end{theorem}

\begin{proof}
By \cite[Theorem~4.3]{L1S} every instance is good. In a graph with no even vertex,
$A=V\setminus N[c]$ is admissible, and avoiding it means $V(Q_c)\subseteq N[c]$. EE is
\cite[Corollary~4.4]{L1S}.
\end{proof}

So (R1), (R2) and (R3) are three faces of one statement. In (R3) the instances that EE needs are not
blocked at the first level; the trouble appears deeper in the reductions. In $\mathrm{CP}(3)$ plus an
apex (the instance where recursive certificates fail \cite{L1S}) every Case~I target of
$(K-uz,z,\{u\})$ lies in $N(c)$, so every step is tight at $c$; of the $456$ good decompositions of
$G-gx$, $320$ transfer successfully and $168$ satisfy the condition of Theorem~\ref{thm:l1tight}
(E26(D)).

\subsection{PTM: walks through the short end's neighbourhood}
In the walk theorem \cite[Theorem~3.2]{TM}, the walk $P=p_0\cdots p_k$ from the long end $y$ to the
short end $a$ is \emph{admissible} if no inner vertex lies in $N(x)$, where $x$ is the endless short
end of the other edge. Walking back is a sequence of Fan transfers with targets $p_i$; admissibility
keeps $x$ out of $Z(p_i)$. If $p_i\in N(x)$, then $Z(p_i)=\{x\}$ and the step is a tight transfer
with $p=x$.

\begin{theorem}[tight walk theorem]\label{thm:tightwalk}
In the walk theorem, drop the condition $p_i\notin N(x)$ $(0<i<k)$. If at each step $i$ with
$p_i\in N(x)$ the current decomposition has no vertex of $\{p_{i+1}\}\cup N(p_i)$ with $x$ as its
predecessor of $p_i$ (for instance, the end on $x$'s side of the path through $xp_i$ lies outside
$\{p_{i+1}\}\cup N(p_i)$), then the conclusion of the walk theorem holds.
\end{theorem}

\begin{proof}
The proof of \cite[Theorem~3.2]{TM} uses admissibility only to apply Fan's lemma at the targets
$p_i$ with $Z(p_i)=\emptyset$. At a step with $Z(p_i)=\{x\}$, use Theorem~\ref{thm:protected} with
$B=\{x\}$ instead. It keeps $b(x)=0$ and gives the same vertex-set conclusion, so the three
separations survive as before.
\end{proof}

This is conditional on the intermediate decompositions, and it is not formalised. In the data it
closes every one of the $68$ instances of order $8$ that have no admissible walk (E26(C)).

\subsection{Summary}
\begin{center}\small
\begin{tabular}{@{}llll@{}}
\toprule
residual & blank $p$ & failure & literally tight?\\
\midrule
$L_\infty$ (R1) & far end $y$ of the leaf path & transfer & yes\\
$L_1^*$ blocked, $u\in N(c)$ (R2) & $c$ & hypothesis and transfer & yes, given the hypothesis\\
$L_1^*$ blocked, $u\notin N(c)$ (R2) & --- & hypothesis (admissibility) & no\\
EE, two common neighbours (R3) & as (R2) & reduces to (R2) & ---\\
PTM, no admissible walk (R4) & short end $x$ & transfer & yes\\
PTM, $K_4$/$W_4$ base (R4) & --- & base separation & no\\
\bottomrule
\end{tabular}
\end{center}

The common form, as a lemma-shaped statement:
\begin{conjecture}[tight transfer]\label{conj:tight}
In a tight transfer (Definition~\ref{def:tight}) arising from Lov\'asz's induction for $L_1^*$ or for
GL, among the decompositions that the induction hypothesis can supply, some has the end of the path
through $pw$ on $p$'s side unreached by the forced chain.
\end{conjecture}
As stated this is a statement about the induction hypothesis, and E26(B) shows that some
strengthening of the hypothesis is needed: no choice of the target alone suffices.

\section{Flowers do not produce Lov\'asz decompositions}\label{sec:flower}

The flower lemma \cite[\S7]{ES} takes paths $P_1,\dots,P_k$ closed into cycles $C(P_i)$ and a
transversal $L$ with one edge $e_i$ of each cycle, and returns the $k+1$ paths $C(P_i)-e_i$ and $L$,
with two ends at every vertex of $L$ and none elsewhere\lean{flower\_decomp}. Receipt-or-flower adds
untouched paths\lean{flower\_decomp\_rest}. ``Equality forces a flower'': if $b(s)=|T|$, $b=1$ on
$T$ and $0$ elsewhere, the paths join $s$ to $T$ and close into cycles\lean{hub\_cycles}.

\begin{proposition}\label{prop:noflower}
\begin{enumerate}
\item If a decomposition satisfies the hypothesis of \texttt{hub\_cycles} and has at most one end at
every vertex, then $|T|\le1$.\lean{hub\_cycles\_card\_le\_one}
\item Every decomposition with the end counts produced by the flower lemma or by
receipt-or-flower has a vertex with at least two ends (every vertex of $L$).
\lean{flower\_decomp\_not\_le\_one, flower\_decomp\_rest\_not\_le\_one}
\end{enumerate}
\end{proposition}

\begin{proof}
(1) $b(s)=|T|\le1$. (2) The first vertex of $L$ has two ends from the flower, plus those of the
untouched paths.
\end{proof}

So no flower move ends in a decomposition with $b\le1$, and in particular never in a Lov\'asz
decomposition; and the state in which equality forces a flower never arises inside Lov\'asz's
induction except trivially (one path from $s$ to the single vertex of $T$). Each of (R1)--(R4)
asks for a Lov\'asz decomposition, so no flower move closes any of them, at any order. The flower
move spends a \emph{spare path} \cite[\S7]{ES}: it is a tool for the rooted
tests (two ends at the root, budget one more path than the parity minimum), not for Lov\'asz-type
invariants, which have no spare path. The tight configuration of Definition~\ref{def:tight} is a
different equality case of Fan's counting: equality at a \emph{target} with no ends, not at a
\emph{hub} carrying $|T|$ ends.

\section{Test 1 at a vertex of degree four}\label{sec:test1}

The motivation for PTM$_2$ at $d(z)=4$ was the first deletion test of the spine through STL
\cite{TM}. That test needs no PTM:

\begin{theorem}\label{thm:test1}
Let $z$ be the only even vertex of $J$, with $d(z)=4$ and $J-z$ connected. Then for every
$r\in N(z)$, $J-z$ has a partition into $(|J|-1)/2$ paths with two ends at $r$.
\lean{test1\_deg\_four}
\end{theorem}

\begin{proof}
$J-z$ has even order $|J|-1$ and exactly four even vertices, the neighbours of $z$ (connectivity
makes every degree positive). By the endpoint matroid for four even vertices \cite{ES}
(\texttt{emc\_four}, from two-step faithful descent, which uses only Lov\'asz's theorem and Fan's
transfers), for any two further even vertices $x,y$ one of $\{r,x\}$, $\{r,y\}$ is realisable. A
realisation of a pair has $\#\mathrm{odd}/2+2=(|J|-1)/2$ paths and two ends at $r$.
\end{proof}

So the $K_4$/$W_4$ base failures of PTM do not obstruct Test~1 at degree four; PTM matters as a
composition mechanism for larger degrees, not for this consequence. (The deletion tests at $d=4$ were
already settled, by published results and by two-step descent; the point here is that the spine's
Test~1 at degree four is now a Lean theorem that does not depend on PTM.)

\section{Computations}\label{sec:comp}

\emph{This section is computational. No proof above depends on it.} Scripts are in
\nolinkurl{manuscripts/tight-fan/scripts/}. \texttt{tf.py} enumerates all path decompositions with
prescribed end counts by transition systems and runs Fan's chain exactly, with every choice;
it agrees with the oracle \texttt{l1s} of \cite{L1S} on all $1{,}772$ random avoidance queries of
orders $5$--$7$ (\texttt{xcheck.py}). In every Case~I step tested below the Fan image was unique,
as Proposition~\ref{prop:forced} predicts.

\textbf{E26(A), $L_1^*$ blocked.} (\texttt{exp1c.py}, \texttt{exp1d.py}, \texttt{exp1e.py}.)
\begin{itemize}
\item Order $9$, the $2{,}414$ blocked instances with $|A|\le2$ that survive every reduction of
\cite{L1S} one level deep: every one has a target $u$ with $(G-gu,c,A)$ good although
inadmissible. In all $162+33$ ``separated'' instances and in $637$ of the $1{,}019$ instances with
$g\in A$ and $d(g)=4$, such a $u$ lies outside $N(c)$. In all $1{,}200$ instances with $g\in A$ and
$d(g)=2$, and in the other $382$ with $d(g)=4$, every good $u$ lies in $N(c)$.
\item With $A=V\setminus(N[c]\cup N(g))$ (when this set equals its boundary $\partial C$),
order $9$, all $12{,}622$ blocked instances: every one has a good target, and in $2{,}938$ of them
every good target lies in $N(c)$.
\item Order $7$, every good decomposition $\cD'$ of $G-gu$ and its forced Fan image. With
$|A|\le2$: among the transfers with $u\in N(c)$, $896$ have a chain avoiding $c$, and all give a
good result; $696$ have a chain through $c$, and all give a bad result; the $56$ transfers with
$u\notin N(c)$ all succeed. With $A$ maximal: $930$ and $515$, and $618$ with $u\notin N(c)$, all
successful.
\item The sufficient condition of Theorem~\ref{thm:l1tight} holds for some good $\cD'$ in every
tight instance of order $7$ ($14$ with $|A|\le2$, $49$ with $A$ maximal), and it never failed to
deliver. ORDER9-E1E
\end{itemize}

\textbf{E26(B), $L_\infty$ residual} (\texttt{linf.py}). All connected graphs of order $7$ and $9$
with a leaf $\ell$ at $c$, exactly one even vertex $z$, $d(c,z)=3$, $N(z)$ at distance $2$, and
$d(z)\ge4$ or $z$ on a triangle: $3$ and $133$ instances. Every one satisfies GL. For every target
$u\in N(z)$ and every decomposition $\cD'$ of $K-zu$ with its $\ell$-path in $B$, we ran the forced
chain. Order $9$: $409{,}416$ transfers; the far end $y$ is a neighbour of $u$ in $206{,}780$ of
them; the chain reaches $y$ in $53{,}964$, and exactly these fail. Every instance has a successful
pair $(u,\cD')$. But in $59$ of the $133$ instances \emph{every} target $u$ admits a failing
$\cD'$; in the other $74$ some $u$ succeeds for every $\cD'$.

\textbf{E26(C), PTM} (\texttt{ptmwalk.py}). The $832$ NOBASE instances of order $8$ of
\cite[E24(F)]{TM}, with walks of any length, admissible or not, base decompositions enumerated
exactly, and the walk back run with exact Fan chains; every success is checked directly on $G$.
\begin{itemize}
\item The $68$ instances with no admissible walk: all $68$ are closed by a tight walk; trying walks
in order of length, the shortest successful walk has length $2$ in $63$ of them and $3$ in $5$.
\item The $764$ $K_4$/$W_4$-type instances: PTM-SAMPLE
\end{itemize}

\textbf{E26(D), the named instances} (\texttt{ee\_inst.py}). $\mathrm{CP}(3)$ plus an apex
(\texttt{F]\char`~vw}) and the double $K_4$ (\texttt{FQhVw}): for each EE pair, the instances
$(K-uz,z,\{u\})$ and the triangle instances, none blocked at the first level. In $\mathrm{CP}(3)$ plus
apex every Case~I target of $(K-uz,z,\{u\})$ is adjacent to $c$ (tight at $c$): $456$ good
decompositions of $G-gx$ per target, $320$ successful transfers, $168$ satisfying the condition of
Theorem~\ref{thm:l1tight}. In the double $K_4$, every target of $(K-uz,z,\{u\})$ is tight at $c$
and all $108$ good decompositions per target satisfy the condition; the triangle instances have their
tight targets with no good decomposition at all (hypothesis side), and their other targets succeed.

\section{What remains: the precise obstruction}\label{sec:open}

\begin{enumerate}
\item \emph{The flower move closes nothing here} (Section~\ref{sec:flower}); this is a structural
fact, not a matter of small cases.
\item \emph{$L_\infty$.} The obstruction $Y$ is: a residual instance together with, for each target
$u\in N(z)$, an induction decomposition of $K-zu$ whose forced chain reaches the far end of the
leaf path. It occurs in $59$ of $133$ residual instances of order $9$. Any proof must therefore
strengthen the induction hypothesis to control the decomposition (for example, by prescribing that
the far end of the leaf path, or the end on its side of the path through $yu$, avoid $N(u)$), not
only choose the target.
\item \emph{$L_1^*$.} Two things are needed: an inductive class that contains the inadmissible but
good instances $(G-gu,c,A)$ (every irreducible blocked instance of order $9$ has one); and, when all
good targets lie in $N(c)$, a decomposition meeting the condition of Theorem~\ref{thm:l1tight}.
\item \emph{PTM.} The no-walk instances are the tight configuration and the tight walk handles
them in the data; the $K_4$/$W_4$ base failures are a base-side obstruction not touched by Fan
equality; Test~1 at $d(z)=4$ no longer depends on them (Theorem~\ref{thm:test1}).
\end{enumerate}

\section{Formalisation}\label{sec:lean}

\texttt{lean/GallaiLean/TightFan.lean} (Lean~4, Mathlib; no \texttt{sorry}; \texttt{\#print
axioms} shows only \texttt{propext}, \texttt{Classical.choice}, \texttt{Quot.sound}; the full
library builds).
\begin{itemize}
\item \texttt{fanPre\_protect}, \texttt{fan\_exists\_chain\_protected},
\texttt{fan\_grouped\_protected\_supp}, \texttt{fan\_single\_protected\_supp}:
Theorem~\ref{thm:protected}.
\item \texttt{fanPred\_unique} (with \texttt{pred\_unique\_of\_isPath},
\texttt{succ\_unique\_of\_isPath}): Proposition~\ref{prop:forced}.
\item \texttt{l1\_tight\_transfer}: Theorem~\ref{thm:l1tight}.
\item \texttt{LInfty}, \texttt{lInfty\_of\_l1StarGen}, \texttt{lInfty\_of\_l1Blocked}:
Theorem~\ref{thm:linfty}.
\item \texttt{hub\_cycles\_card\_le\_one}, \texttt{flower\_decomp\_not\_le\_one},
\texttt{flower\_decomp\_rest\_not\_le\_one}: Proposition~\ref{prop:noflower}.
\item \texttt{test1\_deg\_four}: Theorem~\ref{thm:test1}.
\end{itemize}
Not formalised: Corollary~\ref{cor:oneblank}(1), Proposition~\ref{prop:linfty},
Theorem~\ref{thm:tightwalk}, and the computations.

\begin{thebibliography}{9}
\bibitem{Fan05} G.~Fan, Path decompositions and Gallai's conjecture, J.~Combin.\ Theory Ser.~B 93
(2005) 117--125.
\bibitem{L68} L.~Lov\'asz, On covering of graphs, in: Theory of Graphs (Proc.\ Colloq.\ Tihany,
1966), Academic Press, New York, 1968, 231--236.
\bibitem{ES} J.~Williams, Endpoint structure of path decompositions, manuscript
(\texttt{manuscripts/endpoint-structure}), 2026.
\bibitem{L1S} J.~Williams, The avoidance conjecture $L_1^*$, manuscript
(\texttt{manuscripts/l1star}), 2026.
\bibitem{PE} J.~Williams, Prescribed end edges in Lov\'asz decompositions with one even vertex,
manuscript (\texttt{manuscripts/prescribed-end}), 2026.
\bibitem{RD} J.~Williams, Route descent, manuscript (\texttt{manuscripts/route-descent}), 2026.
\bibitem{TM} J.~Williams, Two terminal edges at once: the terminal-matching property, manuscript
(\texttt{manuscripts/terminal-matching}), 2026.
\end{thebibliography}
\end{document}
